\documentclass[11pt,letterpaper]{article}
\usepackage[T1]{fontenc}
\usepackage{lmodern,amsmath,amssymb,amsthm,microtype}
\usepackage[margin=1in]{geometry}
\usepackage[hidelinks]{hyperref}
\providecommand{\tightlist}{\setlength{\itemsep}{0pt}\setlength{\parskip}{0pt}}
\DeclareUnicodeCharacter{2264}{\ensuremath{\le}}
\DeclareUnicodeCharacter{2265}{\ensuremath{\ge}}
\DeclareUnicodeCharacter{2260}{\ensuremath{\ne}}
\DeclareUnicodeCharacter{2212}{\ensuremath{-}}

\DeclareUnicodeCharacter{2211}{\ensuremath{\sum}}
\title{Block deficits in the five-isolate height-six Frankl case}
\author{Aafi Khan Malek\\ \small Independent researcher, Ahmedabad, India\\
\small\href{mailto:aafimalek2032@gmail.com}{aafimalek2032@gmail.com}}
\date{28 November 2026}
\begin{document}
\maketitle
We use the setup of the
\href{https://github.com/Aafimalek/novel-discoveries/tree/2d50235bd73e18b1c1371f96dd7fc3540514e0ee/frankl-height-six/verification/at-least-four-isolates}{preceding paper}.
A weak family is finite, nonempty, union-closed, contains no empty
member and no strict chain of seven members, and has every point
in at most half of its members. Ground minimality is taken among
such weak families. A proper forcing edge means that a distinct
second point appears whenever the first point appears.
\textbf{Written theorem with a finite integer replay. Not yet checked
end to end in Lean.} Let a weak minimum-ground counterexample at height
at most six have exactly five isolated points and at least two
forcing-star centers. Its omitted-star pair design has maximum
block-degree deficit at most two. No unary-star hypothesis is needed for
this theorem. An additional argument below bounds the number of
deficit-two centers without an extra hypothesis. General cases with
deficits zero, one, and two, and the single-source branch, are not
excluded.

\section*{Setup}\label{setup}

Let k count sources, Q = binom(k,2), N count omitted-star blocks, and
g\_v count blocks through v. Write d\_v = k-1-g\_v and d = max d\_v.
Choose a minimum-degree center p, with g = g\_p = k-1-d. At any center v
the other points of its blocks partition the other k-1 centers, so

\begin{verbatim}
d_v = sum over B containing v of (|B|-2),   1 <= g_v <= k-1.
\end{verbatim}

Let H(m) = 1+m+binom(m,2). For distinct anchors p,q and anchor block B,
omit all of B. The exterior blocks through p give g-1 groups of equal
occurrence signature. The sharp compressed Sauer bound on the
pair-avoidance fiber is H(g+4): the g-1 groups plus five isolates. The
fiber\textquotesingle s lower bound is

\begin{verbatim}
|F(p,q)| >= n+k+N-g_p-g_q >= N+k+7+d+d_q.       (1)
\end{verbatim}

Here n>=2k+5, and exactly N-g\_p-g\_q+1 blocks avoid the anchor pair. If
q\textquotesingle s exterior groups merge r independent pairs of the p
groups, the upper bound lowers to H(g+4-r).

For any block of size r>=2 its loss of blocks from the all-pair design is
binom(r,2)-1 <= 2r(r-2)/3. Thus

\begin{verbatim}
Q-N <= (2/3) sum_v d_v <= 2kd/3.                (2)
\end{verbatim}

We also use the preceding incidence Cauchy bound N >= k \(g^2\)/(k-1+g).

\section{Remove deficits at least six and restrict all
others}\label{1-remove-deficits-at-least-six-and-restrict-all-others}

Drop d\_q in (1). Comparing the starting bounds, substituting the Cauchy
bound, and clearing the positive denominator yields

\begin{verbatim}
(d-12)g^2+(d-12)g+d(4d-6) <= 0.
\end{verbatim}

For d>=12 the left side is positive, so d<=11. Substituting (2) instead
shows that the slack between the upper and lower bounds is at most

\begin{verbatim}
Delta = (3-d/3)k + (d^2-9d)/2.                     (3)
\end{verbatim}

The first extra merger costs g+4 = k+3-d. For d>=6, the difference
\(\Delta\)-(k+3-d) is nonincreasing in k and at k=d+2 equals \((d^2-13d+6)/6\),
negative for 6<=d<=11. Therefore no extra merger is possible under those
deficits.

For d=3,4,5, two extra independent mergers would cost 2(k+3-d)-1
\textgreater{} \(\Delta\). Thus only one extra merger is possible. In either
reasoning, take any block C not through p and any \(q\in C\). Its other points
are outside B, since distinct blocks share at most one point. They lie
in distinct p groups, since otherwise C and a block through p would
share at least two points. Consequently C forces \textbar C\textbar-2
mergers. For d>=6, all nonpair blocks must pass through p; for d=3,4,5,
all blocks outside p have size two or three.

When all nonpair blocks pass through p, their positive excesses
e=\textbar B\textbar-2 sum to d. The lost-block count Q-N = sum
\((e^2+3e)/2\) is at most \((d^2+3d)/2\). With this bound instead of (2), (3)
improves to slack at most (d-3)(d-k). This excludes every
d\textgreater3; for d=3 the slack is at most zero, and choosing q in any
nonpair block through p supplies its d\_q>=1, also impossible. We are
left with d=3,4,5 and at least one off-p triple.

\section{Group charging bounds the number of
centers}\label{2-group-charging-bounds-the-number-of-centers}

The nontrivial groups through p have sizes t\_i=e\_i+1, where the
positive e\_i partition d. Let m=sum\_i t\_i. The other k-1-m groups are
singletons. For any q in a singleton group, if it lies on s>=2 off-p
triples, their pairs of other vertices must all join the same two p
groups. Otherwise two distinct links between group labels give two
independent mergers. The s pairs are vertex-disjoint because two
distinct triples sharing q cannot share a second point. In particular
the two partner groups are nontrivial and s<=min(t\_i,t\_j). Each pair
can be used by at most one q, by pair uniqueness.

Give each singleton q baseline deficit one, then charge s-1 when s>=2.
For a fixed pair of nontrivial groups, the sum of these charges over
assigned singleton vertices is at most

\begin{verbatim}
(1-1/min(t_i,t_j)) t_i t_j = t_i t_j-max(t_i,t_j).
\end{verbatim}

Every vertex in the m nontrivial groups, and p itself, has deficit at
most d. Hence

\begin{verbatim}
sum_v d_v <= k + C,
C = (d-1)(1+m) + sum_{i<j}(t_i t_j-max(t_i,t_j)).
\end{verbatim}

The maxima over partitions of d=3,4,5 are 20,39,64. This tiny
computation is checked in the replay. Substituting this refined deficit
bound into (2), the slack in the starting comparison improves to

\begin{verbatim}
(3-d+2/3) k + (d^2-9d)/2 + 2C/3.
\end{verbatim}

Now anchor at any q in an off-p triple. This costs one extra merger
k+3-d and at least one from d\_q in (1). For d=3,4,5 the resulting
inequalities give k<=10,12,14 respectively.

\section{Exact partition check eliminates all those
cases}\label{3-exact-partition-check-eliminates-all-those-cases}

For a fixed partition of d and a fixed k, include all singleton p groups
by setting their e\_i=0. As noted above, each off-p triple takes one
vertex from each of three distinct p groups. For any q in group i, even
a nonsingleton group, all its off-p triples must use the same two
partner groups: the anchor block for p,q drops group i, and two distinct
links on the other labels would force two independent mergers. The
number of its off-p triples cannot exceed d-e\_i or the second-largest
size among the other groups. The latter is because the two selected
partner groups must each supply a different vertex per triple. Therefore
for t off-p triples,

\begin{verbatim}
3t <= sum_i (e_i+1) min(d-e_i, second_largest_other_size_i).
\end{verbatim}

With fewer than three groups there cannot be an off-p triple. The exact
loss is

\begin{verbatim}
Q-N = sum_i (binom(e_i+2,2)-1) + 2t.
\end{verbatim}

In contrast, using the tightened anchor upper H(g+3) in (1) with d\_q>=1
requires

\begin{verbatim}
Q-N >= (d-2)k+(6+7d-d^2)/2+1.
\end{verbatim}

The self-contained standard-Python replay considers every positive
partition of d, every k from d+2 through its proved bound, and all
forced singleton groups. There are 12,24,39 eligible rows with at least
three groups for d=3,4,5. The largest values of the maximum allowable
loss minus the necessary loss are -2, -9, -12, respectively. No row
survives. This completes the maximum-deficit-two conclusion.

\subsection{Optional sharpening in the deficit-two
branch}\label{optional-sharpening-in-the-deficit-two-branch}

Assume now d=2. Let D count the deficit-two centers and M=Q-N. Only
triples and four-blocks can be nonpairs, and the latter do not meet any
other nonpair block since all their vertices are already at deficit two.
Equation (2) gives M<=2(k+D)/3<=4k/3.

For any fixed deficit-two center r, at most four other deficit-two
centers share a nonpair block with r. For any other deficit-two center q
the anchor block is an ordinary pair. Their four exterior identification
units are independent except in the alternating four-cycle of two
mate-pairs on each side, as in the four-isolate theorem. For fixed r
there is at most one such cycle partner q: in the dual simple graph on
triples, where a degree-two center is an edge, the two endpoints of
r\textquotesingle s edge have at most two common neighbors (the dual
maximum degree is three). The cycle partner would have to be the edge
between those two neighbors. No such cycle is possible for r belonging
to a four-block. Call a pair with all four independent units a good
pair.

If any good pair exists, its fiber projects onto at most k-1 coordinates
(k-2 outside centers, five isolates, and four identifications). By (1),
M>=k+10 is necessary, so D>=k/2+15 by the bound on M. Also k>=30 follows
from k+10<=M<=4k/3. This forces all stars to be unary in the good-pair
case: the checked multileaf bound with a second source would give
n+k<=2+5+binom(5,2)=17. Consequently every deficit-two center has at
least two good partners. For every good pair the slack at H(k-1)=Q+1 is
M-k-10<=k/3-10, smaller than k-2. The preceding rank-two atom bound now
forces exactly k-1 supported quotient coordinates with all singleton
omissions. In particular each isolate can be omitted from the full pair
maximum on its own. Taking two good partners for a center r and unioning
the relative omission sets yields a ground complement consisting of
r\textquotesingle s whole star and any specified isolate. For one good
anchor pair the remaining D-2 deficit-two stars therefore give 5(D-2)
extra mixed omission witnesses, distinct from top, single-point
avoidance maxima, source witnesses and pair maxima by the same argument
as in the four-isolate proof. This number exceeds the slack M-k-10,
contradiction.

Thus there is no good pair. If D>=7, the at-most-five exceptions for a
fixed r would give a good pair, so D<=6. This also shows that M>=k+10 is
impossible: that inequality alone forces D>=k/2+15 and hence a good pair.
The low-loss branch with D<=6 is not ruled out.

\section*{Verification map}\label{verification-map}

Run \texttt{python3 deficit\_replay.py} for the finite integer check.
The formal prerequisites are the checked pair-design, link-signature,
compressed Sauer, witness-count and multileaf statements from the earlier
development. The companion note gives the precise Lean declaration names.
The supplementary low-loss argument also uses the earlier written atom
bound and mixed-omission argument. Its new results are not Lean theorems.

\paragraph{Contributions.}
Aafi Khan Malek selected and directed the research. ChatGPT assisted
with the proof, verification and writing. Codex was not used.
The note has not been peer reviewed.
\end{document}
